\chapter*{\centering Abstract}
\addcontentsline{toc}{chapter}{ABSTRACT}

\hs Health insurance underwriting in emerging markets such as Cambodia faces structural challenges including low insurance penetration, fragmented healthcare infrastructure, and limited longitudinal health records. Traditional static rule-based underwriting systems are suboptimal because they ignore feature interactions, cannot adapt to portfolio drift, and may inadvertently discriminate against demographic segments. This thesis investigates whether contextual bandits --- a class of online learning algorithms for sequential decision-making under uncertainty --- can replace static rules in emerging-market health insurance underwriting.

The research designs and implements a contextual bandit framework with three algorithms (LinUCB, LinTS, and Epsilon-Greedy) and compares them against a static XGBoost rule baseline on a synthetic dataset of 2{,}000 Cambodian health insurance applicants anchored on the Cambodia Demographic and Health Survey 2021--22 \citep{NIS2023}. A profit-based actuarial reward simulator evaluates four underwriting actions (standard, rated, decline, refer); Population Stability Index (PSI) sliding-window guardrails and the U.S.\ EEOC four-fifths rule jointly monitor regional and occupational fairness. All headline results are reported as means across 20 independent seeds with paired Wilcoxon signed-rank tests and bootstrap 95\% confidence intervals; EXP-008 (human-in-the-loop, §5.4) is reported at SEED = 42 pending a 20-seed rerun.

Four controlled experiments validate the framework. EXP-005 shows that LinUCB achieves \$90{,}540 cumulative reward over 5{,}000 rounds against the static baseline's \$72{,}292 ($+25\%$, Wilcoxon $p < 0.001$, Cohen's $d = 2.98$), and cuts average regret in the final 500 rounds from \$9.01 to \$2.20 per round. EXP-006 confirms that the maximum sliding-window PSI remains in the GREEN/AMBER band for both region (0.082) and occupation (0.123) while the EEOC four-fifths rule is satisfied with parity ratios of 85.7\% and 90.1\% respectively. EXP-007 ranks the algorithms by cumulative regret, with LinTS and LinUCB statistically indistinguishable and both decisively outperforming the static and uniform-exploration baselines. EXP-008 wraps LinUCB in a human-in-the-loop layer and raises cumulative reward by 6.5\% over the mathematical-REFER baseline at a human-review cost of 2.6\% of reward. A post-hoc robustness analysis identifies an inadmissible constant policy as a theoretical performance ceiling, bounding the interpretation of this result.

The results provide a rigorous proof of concept that contextual bandits can outperform the static rule-based underwriting they would replace on profitability while maintaining demographic fairness --- while remaining computationally lightweight enough for deployment on low-resource mobile infrastructure; the inadmissible constant ceiling identified above bounds, but does not overturn, this contribution against admissible alternatives. The thesis contributes a reproducible 20-seed experimental harness, a Cambodia-calibrated synthetic dataset, sliding-window PSI and EEOC four-fifths fairness monitoring, and a human-in-the-loop wrapper that may inform future research and industry practice in emerging-market algorithmic insurance.

\bigskip
\noindent\textbf{Keywords:} Contextual Bandits, LinUCB, Thompson Sampling, Reinforcement Learning, Health Insurance Underwriting, Cambodia, Algorithmic Fairness, Population Stability Index, Actuarial Science, Online Learning
